% !TeX program = xelatex
% !TeX encoding = UTF-8
% =========================================================
% hustreport.cls - Complete Source Code Analysis Tutorial
% Version: v3.0
% Language: English
% =========================================================
\documentclass[linux]{hustreport}

% Document Information
\title{hustreport.cls Source Code Analysis}
\author{HUST Report Development Team}

\begin{document}

% =========================================================
% Title Page
% =========================================================
\begin{titlepage}
    \centering
    \vspace*{3cm}
    {\Huge\bfseries\heiti hustreport.cls\par}
    \vspace{0.5cm}
    {\Large\songti HUST University Report Template Class\par}
    \vspace{1cm}
    {\large Complete Source Code Analysis Tutorial\par}
    \vspace{2cm}
    {\large Version 3.0\par}
    \vspace{0.5cm}
    {\large \today\par}
    \vfill
    {\small This document provides a comprehensive module-by-module analysis of the hustreport.cls class file.\par}
\end{titlepage}

\tableofcontents
\newpage

% =========================================================
% Chapter 1: Introduction and Overview
% =========================================================
\section{Introduction and Overview}

\subsection{Template Purpose}

The \texttt{hustreport.cls} is a \LaTeX{} document class designed for Heilongjiang University of Science and Technology (HUST) undergraduate students. It provides a complete solution for:

\begin{itemize}
    \item Social investigation reports
    \item Internship reports
    \item Graduation theses (with \texttt{thesis} option)
\end{itemize}

\subsection{Architecture Overview}

The class file consists of \textbf{20 functional modules} spanning approximately 1445 lines of code:

\begin{enumerate}
    \item \textbf{Module 0-1}: Version declaration and class options
    \item \textbf{Module 2}: Package loading (categorized by function)
    \item \textbf{Module 3}: User information interface
    \item \textbf{Module 4}: Smart font configuration
    \item \textbf{Module 5-6}: Page layout and section titles
    \item \textbf{Module 7-8}: Figure/table formatting and PDF properties
    \item \textbf{Module 9-10}: Special environments and cover generation
    \item \textbf{Module 11-20}: Advanced features (code blocks, review system, etc.)
\end{enumerate}

\subsection{Compilation Requirements}

This template \textbf{must} be compiled with XeLaTeX:

\begin{verbatim}
xelatex your_document.tex
bibtex your_document      % If using bibliography
xelatex your_document.tex
xelatex your_document.tex
\end{verbatim}

% =========================================================
% Chapter 2: Module 0-1 - Version and Class Options
% =========================================================
\newpage
\section{Module 0-1: Version Declaration and Class Options}

\subsection{Version Declaration (Lines 1-14)}

Every \LaTeX{} class file must declare its format requirements and provide identification:

\begin{verbatim}
% Line 12-13: Format and class identification
\NeedsTeXFormat{LaTeX2e}
\ProvidesClass{hustreport}[2026/01/05 v3.0 HUST Report Class]
\end{verbatim}

\textbf{Key Points:}
\begin{itemize}
    \item \verb|\NeedsTeXFormat{LaTeX2e}| ensures compatibility with modern \LaTeX{}
    \item \verb|\ProvidesClass| provides version information displayed in logs
\end{itemize}

\subsection{Boolean Switches (Lines 19-24)}

The class uses boolean switches to control conditional compilation:

\begin{verbatim}
% Platform detection switches
\newif\ifhust@oswin    \hust@oswinfalse
\newif\ifhust@osmac    \hust@osmacfalse
\newif\ifhust@oslinux  \hust@oslinuxfalse
\newif\ifhust@optfonts \hust@optfontsfalse
\newif\ifhust@thesis   \hust@thesisfalse
\end{verbatim}

\textbf{Implementation Pattern:}
\begin{itemize}
    \item \verb|\newif| creates a boolean toggle with default \texttt{false}
    \item Naming convention: \verb|\ifhust@<name>| uses internal prefix \verb|hust@|
    \item All switches initialize to \texttt{false}, enabling opt-in behavior
\end{itemize}

\subsection{Option Declaration (Lines 27-31)}

Options are declared using \verb|\DeclareOption|:

\begin{verbatim}
% Platform options
\DeclareOption{win}{\\hust@oswintrue}
\DeclareOption{mac}{\\hust@osmactrue}
\DeclareOption{linux}{\\hust@oslinuxtrue}
\DeclareOption{fonts}{\\hust@optfontstrue}
\DeclareOption{thesis}{\\hust@thesistrue}
\end{verbatim}

\subsection{Option Forwarding and Processing (Lines 33-35)}

Unrecognized options are forwarded to the base class:

\begin{verbatim}
\DeclareOption*{\PassOptionsToClass{\CurrentOption}{ctexrep}}
\ProcessOptions\relax
\end{verbatim}

\subsection{Base Class Loading (Lines 37-43)}

The template is built on \texttt{ctexrep} for Chinese support:

\begin{verbatim}
\LoadClass[a4paper,UTF8,zihao=-4,linespread=1.5]{ctexrep}
\end{verbatim}

\textbf{Parameters Explained:}
\begin{optionlist}
    \option{a4paper}{Standard A4 paper size (210mm × 297mm)}
    \option{UTF8}{UTF-8 encoding for Chinese characters}
    \option{zihao=-4}{Default font size: xiaosi (12pt equivalent)}
    \option{linespread=1.5}{1.5x line spacing for readability}
\end{optionlist}

% =========================================================
% Chapter 3: Module 2 - Package Loading
% =========================================================
\newpage
\section{Module 2: Package Loading}

\subsection{Organization Philosophy}

Packages are organized by functional category for maintainability:

\subsection{Page and Layout Packages (Lines 51-55)}

\begin{verbatim}
\RequirePackage{geometry}   % Page margins
\RequirePackage{fancyhdr}   % Headers and footers
\RequirePackage{float}      % Float placement control
\RequirePackage[sort&compress,numbers,square]{natbib}
\end{verbatim}

\subsection{Title and TOC Packages (Lines 57-59)}

\begin{verbatim}
\RequirePackage{titletoc}   % TOC formatting
\RequirePackage{titlesec}   % Section title formatting
\end{verbatim}

\subsection{Mathematics and Physics (Lines 61-66)}

\begin{verbatim}
\RequirePackage{amsmath,amssymb,amsfonts,amsthm}
\RequirePackage{physics}    % Physics notation shortcuts
\RequirePackage{siunitx}    % SI units
\RequirePackage{newtxmath}  % Math font optimization
\RequirePackage{calc}       % Length calculations
\end{verbatim}

\subsection{Graphics and Color (Lines 68-72)}

\begin{verbatim}
\RequirePackage{graphicx}   % Image inclusion
\RequirePackage{subcaption} % Subfigures
\RequirePackage{tikz}       % Drawing and positioning
\RequirePackage[dvipsnames,svgnames,x11names]{xcolor}
\end{verbatim}

\subsection{Tables and Lists (Lines 74-84)}

\begin{verbatim}
\RequirePackage{booktabs}   % Professional tables
\RequirePackage{threeparttable}
\RequirePackage{caption}
\RequirePackage{tabularx}   % Auto-width tables
\RequirePackage{xltabular}  % Long tables with page breaks
\RequirePackage{multirow}   % Multi-row cells
\RequirePackage{makecell}   % Cell line breaks
\RequirePackage{array}      % Column format extensions
\RequirePackage{enumitem}   % List customization
\RequirePackage{diagbox}    % Diagonal box headers
\end{verbatim}

% =========================================================
% Chapter 4: Module 3 - Information Interface
% =========================================================
\newpage
\section{Module 3: Information Interface}

\subsection{Internal Storage Variables (Lines 97-102)}

The class defines internal variables using \verb|@| prefix:

\begin{verbatim}
\makeatletter
\newcommand{\@major}{}      % Major/Program
\newcommand{\@dept}{}       % Department
\newcommand{\@advisor}{}    % Advisor
\newcommand{\@team}{}       % Team members
\newcommand{\@studentID}{}  % Student ID
\newcommand{\@reportID}{}   % Report number
\end{verbatim}

\subsection{Public Interface Commands (Lines 105-110)}

User-facing commands that set the internal variables:

\begin{verbatim}
\newcommand{\major}[1]{\def\@major{#1}}
\newcommand{\dept}[1]{\def\@dept{#1}}
\newcommand{\advisor}[1]{\def\@advisor{#1}}
\newcommand{\team}[1]{\def\@team{#1}}
\newcommand{\studentID}[1]{\def\@studentID{#1}}
\newcommand{\reportID}[1]{\def\@reportID{#1}}
\end{verbatim}

\subsection{Task Book Variables (Lines 113-116)}

Predefined variables for task book content:

\begin{verbatim}
\newcommand{\taskTargetText}{Content pending}
\newcommand{\taskContentText}{Content pending}
\newcommand{\taskExpectedText}{Content pending}
\newcommand{\taskScheduleText}{Schedule pending}
\end{verbatim}

\subsection{Smart Underline Component (Lines 119-123)}

Adaptive underline for form fields:

\begin{verbatim}
\newcommand{\ulineinfo}[2][280pt]{%
    \begin{minipage}[t]{#1}%
        \centering\zihao{3}\songti\uline{#2\hfill}%
    \end{minipage}%
}
\end{verbatim}

% =========================================================
% Chapter 5: Module 4 - Font Configuration
% =========================================================
\newpage
\section{Module 4: Smart Font Configuration}

\subsection{Design Philosophy}

The font system implements a \textbf{fallback chain} with three priority levels:

\begin{enumerate}
    \item \textbf{Highest}: Local \texttt{fonts/} directory
    \item \textbf{Medium}: Platform-specific system fonts
    \item \textbf{Fallback}: Fandol open-source fonts
\end{enumerate}

\subsection{Font Reset (Lines 129-131)}

Clear existing font definitions to prevent conflicts:

\begin{verbatim}
\let\songti\relax
\let\heiti\relax
\let\kaiti\relax
\end{verbatim}

\subsection{English Font Configuration (Lines 147-167)}

\begin{verbatim}
\IfFileExists{fonts/times.ttf}{
    % Local Times font found
    \setmainfont{times.ttf}[
        Path = fonts/,
        BoldFont = timesbd.ttf,
        ItalicFont = timesi.ttf,
        BoldItalicFont = timesbi.ttf
    ]
}{
    % System fallback
    \ifhust@oswin
        \setmainfont{Times New Roman}
    \else
        \ifhust@osmac
            \setmainfont{Times New Roman}
        \else
            \setmainfont{TeX Gyre Termes}
        \fi
    \fi
}
\end{verbatim}

\subsection{Chinese Font Configuration (Lines 171-254)}

The Chinese font setup uses conditional checks for each font file:

\begin{verbatim}
\IfFileExists{fonts/SimSun.ttf}{
    % Local Chinese fonts available
    \setCJKmainfont{SimSun.ttf}[
        Path = fonts/,
        BoldFont = \hust@cjk@bold,
        ItalicFont = \hust@cjk@italic,
        AutoFakeBold = \hust@cjk@fakebold
    ]
    % ... additional font family definitions
}{
    % Platform-specific system fonts
    \ifhust@oswin
        % Windows: SimSun, SimHei, KaiTi
    \else\ifhust@osmac
        % macOS: Songti SC, Heiti SC, Kaiti SC
    \else
        % Linux: Fandol fonts
    \fi\fi
}
\end{verbatim}

% =========================================================
% Chapter 6: Module 5-6 - Page Layout and Section Titles
% =========================================================
\newpage
\section{Module 5-6: Page Layout and Section Titles}

\subsection{Conditional Page Layout (Lines 260-301)}

The layout switches based on the \texttt{thesis} option:

\begin{verbatim}
\ifhust@thesis
    % Thesis mode: special header with logo
    \geometry{
        a4paper,
        top=3.7cm, bottom=2.5cm,
        left=3.9cm, right=2.5cm,
        headheight=3cm, headsep=0.5cm
    }
    
    % Custom header with school logo
    \pagestyle{fancy}
    \fancyhf{}
    \renewcommand{\headrulewidth}{2.5pt}
    \fancyhead[L]{\includegraphics[height=2.6cm]{logo.png}}
    \fancyhead[C]{...school name and title...}
    \fancyfoot[C]{\thepage}
\else
    % Standard report mode
    \geometry{left=3.017cm, right=3.017cm, 
              top=2.54cm, bottom=2.54cm}
    \pagestyle{plain}
\fi
\end{verbatim}

\subsection{Counter Configuration (Lines 305-313)}

Section numbering uses Arabic numerals with chapter binding:

\begin{verbatim}
\renewcommand{\thesection}{\arabic{section}}
\renewcommand{\thesubsection}{\thesection.\arabic{subsection}}
\renewcommand{\thesubsubsection}{\thesubsection.\arabic{subsubsection}}

% Bind counters to sections (reset per chapter)
\counterwithin{equation}{section}
\counterwithin{figure}{section}
\counterwithin{table}{section}
\end{verbatim}

\subsection{Section Title Styles (Lines 325-352)}

Using \texttt{ctexset} for Chinese-aware formatting:

\begin{verbatim}
\ctexset{
    section = {
        number = \chinese{section},
        name = {Chapter ,},
        format = \centering\zihao{3}\heiti\bfseries,
        aftername = \quad,
        beforeskip = 15pt, afterskip = 20pt
    },
    subsection = {
        format = \zihao{4}\heiti\bfseries,
        aftername = \hspace{1em}
    },
    subsubsection = {
        format = \zihao{-4}\songti,
        aftername = \hspace{1em}
    }
}
\end{verbatim}

\subsection{TOC Styles (Lines 358-376)}

Custom table of contents formatting:

\begin{verbatim}
\titlecontents{section}[0pt]
{\addvspace{6pt}\bfseries\heiti\zihao{4}}
{\thecontentslabel\quad}
{}
{\titlerule*[0.5pc]{.}\contentspage}
\end{verbatim}

% =========================================================
% Chapter 7: Module 7-8 - Figures, Tables, and PDF
% =========================================================
\newpage
\section{Module 7-8: Figures, Tables, and PDF Properties}

\subsection{Float Parameters Optimization (Lines 386-392)}

Solving the ``figure jumping'' problem:

\begin{verbatim}
\renewcommand{\topfraction}{0.85}    % Top floats: 85%
\renewcommand{\bottomfraction}{0.70} % Bottom floats: 70%
\renewcommand{\textfraction}{0.10}   % Min text: 10%
\renewcommand{\floatpagefraction}{0.50}
\end{verbatim}

\subsection{Caption Styling (Lines 395-410)}

Professional caption formatting:

\begin{verbatim}
\DeclareCaptionFont{heiti}{\heiti}
\DeclareCaptionFont{fivehao}{\zihao{5}}
\captionsetup{
    format=hang,
    labelsep=quad,
    font={fivehao,heiti},
    skip=6pt
}
\end{verbatim}

\subsection{Numbering Formats (Lines 415-423)}

Chapter-relative numbering:

\begin{verbatim}
\renewcommand{\theequation}{\thesection-\arabic{equation}}
\renewcommand{\thefigure}{\thesection-\arabic{figure}}
\renewcommand{\thetable}{\thesection-\arabic{table}}

% Auto-prefix for references
\labelformat{figure}{Figure~#1}
\labelformat{table}{Table~#1}
\labelformat{equation}{Eq.~(#1)}
\end{verbatim}

\subsection{PDF Properties (Lines 442-448)}

Hyperref configuration:

\begin{verbatim}
\hypersetup{
    hidelinks,
    colorlinks=false,
    bookmarksnumbered=true,
    bookmarksopen=true,
    pdfstartview=FitH
}
\end{verbatim}

% =========================================================
% Chapter 8: Module 9-10 - Special Environments
% =========================================================
\newpage
\section{Module 9-10: Special Environments and Cover}

\subsection{Chinese Abstract Environment (Lines 455-468)}

\begin{verbatim}
\newenvironment{cnabstract}[1]{
    \def\@keyword{#1}
    \clearpage
    \phantomsection
    \addcontentsline{toc}{section}{Abstract}
    \begin{center} \heiti\zihao{3} \@title \end{center}
    \vspace{1ex}
    \begin{center} \heiti\zihao{4} Abstract \end{center}
    \par
}{
    \par\vspace{2ex}
    \noindent{\heiti\zihao{-4} Keywords: }\@keyword
    \newpage
}
\end{verbatim}

\subsection{Bibliography Automation (Lines 471-498)}

\begin{verbatim}
\newcommand{\MyBibliography}[1]{
    \clearpage
    \phantomsection
    \addcontentsline{toc}{section}{References}
    \renewcommand{\section}[2]{}% Suppress auto-heading
    \bibliographystyle{plain}
    \bibliography{#1}
}
\end{verbatim}

\subsection{Appendix Environment (Lines 501-527)}

\begin{verbatim}
\newenvironment{appendixblock}{
    \clearpage
    \setcounter{section}{0}
    \renewcommand{\thesection}{\Alph{section}}
    \ctexset{
        section = {
            name = {Appendix ,},
            format = \centering\zihao{4}\heiti\bfseries
        }
    }
}{}
\end{verbatim}

\subsection{Cover Generation with TikZ (Lines 641-680)}

The cover uses TikZ for absolute positioning:

\begin{verbatim}
\newcommand{\makecover}{
    \begin{titlepage}
        \begin{tikzpicture}[remember picture, overlay]
            % Logo positioning
            \node[anchor=north east, xshift=-2cm, yshift=-2cm] 
                at (current page.north east) {
                \includegraphics[width=0.25\textwidth]{logo.png}
            };
        \end{tikzpicture}
        % Title and info table...
    \end{titlepage}
}
\end{verbatim}

% =========================================================
% Chapter 9: Module 11-12 - Advanced Layouts and Code
% =========================================================
\newpage
\section{Module 11-12: Advanced Layouts and Code Environments}

\subsection{Text-Figure Mixed Layout (Lines 693-737)}

The \texttt{textfigure} environment uses \texttt{xparse}:

\begin{verbatim}
\NewDocumentEnvironment{textfigure}{O{right} m m O{0.4}}{%
    \def\tf@pos{#1}%    % Position: left or right
    \def\tf@path{#2}%   % Image path
    \def\tf@cap{#3}%    % Caption
    \setlength{\tf@imagew}{#4\textwidth}%
    % Calculate text width
    \setlength{\tf@textw}{\textwidth}%
    \addtolength{\tf@textw}{-\tf@imagew}%
    \addtolength{\tf@textw}{-1.5em}%
    % Begin text capture...
}{%
    % Render image + text side by side
}
\end{verbatim}

\subsection{Parallel Figures Environment (Lines 739-761)}

The \texttt{parallelfigures} environment uses \texttt{subfigure} for side-by-side images with automatic (a), (b), (c) numbering:

\begin{verbatim}
\NewDocumentEnvironment{parallelfigures}{m}{%
    \renewcommand{\pf@maincap}{#1}%
    \begin{figure}[H]
        \centering
        \newcommand{\addfig}[3][0.48]{%
            \begin{subfigure}[b]{##1\textwidth} 
                \centering
                \includegraphics[width=\linewidth]{##2}
                \caption{##3}
            \end{subfigure}%
            \hspace{\fill}
        }%
    }{%
        \caption{\pf@maincap}%
    \end{figure}%
}
\end{verbatim}

\textbf{Subfigure Numbering}:
\begin{itemize}
    \item Main figure: Chapter-relative numbering (e.g., Figure 9-1)
    \item Subfigures: Automatic (a), (b), (c) format
    \item Defined in line 415: \verb|\renewcommand\thesubfigure{(\alph{subfigure})}|
\end{itemize}

\subsection{New Multi-Figure Environments (v3.0)}

Two new environments for flexible layouts:

\subsubsection{figurerow - Auto-width Row Layout}

\begin{verbatim}
\begin{figurerow}{Main Title}[columns]
    \figitem{path}{subcaption}
    \figitem{path}{subcaption}
\end{figurerow}
\end{verbatim}

\subsubsection{figuregrid - Grid Layout}

\begin{verbatim}
\begin{figuregrid}{Main Title}
    \gridrow{
        \gridfig[width]{path}{subcaption}
        \gridfig[width]{path}{subcaption}
    }
\end{figuregrid}
\end{verbatim}

\subsection{Code Environment Design (Lines 772-811)}

Color definitions for syntax highlighting:

\begin{verbatim}
\definecolor{hust@codeback}{RGB}{250,250,250}
\definecolor{hust@keyword}{RGB}{0,92,197}
\definecolor{hust@comment}{RGB}{80,161,79}
\definecolor{hust@string}{RGB}{227,17,108}

\lstdefinestyle{huststyle}{
    basicstyle=\small\ttfamily,
    keywordstyle=\bfseries\color{hust@keyword},
    commentstyle=\itshape\color{hust@comment},
    stringstyle=\color{hust@string},
    numbers=left,
    breaklines=true,
    keepspaces=true
}
\end{verbatim}

\subsection{tcolorbox Integration (Lines 814-855)}

Mac-style code block with traffic lights:

\begin{verbatim}
\tcbset{
    hustcodestyle/.style={
        enhanced, breakable,
        colback=hust@codeback,
        title={\textbf{Code \thehustcode:}#1},
        % Mac-style decoration
        overlay unbroken and first={
            \draw[fill=red!70!white] 
                ([xshift=6mm,yshift=-1.5mm]frame.north west) 
                circle (2.5pt);
            \draw[fill=yellow!70!white] 
                ([xshift=11mm,yshift=-1.5mm]frame.north west) 
                circle (2.5pt);
            \draw[fill=green!70!white] 
                ([xshift=16mm,yshift=-1.5mm]frame.north west) 
                circle (2.5pt);
        }
    }
}
\end{verbatim}

% =========================================================
% Chapter 10: Module 13-16 - Abstract, Algorithm, Notice
% =========================================================
\newpage
\section{Module 13-16: Abstract, Algorithm, and Notice Environments}

\subsection{English Abstract (Lines 872-887)}

\begin{verbatim}
\newenvironment{enabstract}[1]{
    \def\@enkeyword{#1}
    \clearpage
    \addcontentsline{toc}{section}{Abstract}
    \begin{center}\bfseries\zihao{4} Abstract \end{center}
    \linespread{1.5}\selectfont
    \zihao{-4}
}{
    \par\vspace{2ex}
    \noindent\textbf{Keywords: } \@enkeyword
    \newpage
}
\end{verbatim}

\subsection{Algorithm Localization (Lines 892-904)}

\begin{verbatim}
\RequirePackage{algorithm}
\RequirePackage{algorithmicx}
\RequirePackage{algpseudocode}

\floatname{algorithm}{Algorithm}
\renewcommand{\algorithmicrequire}{\textbf{Input:}}
\renewcommand{\algorithmicensure}{\textbf{Output:}}
\algrenewcommand\algorithmicindent{2.0em}
\end{verbatim}

\subsection{Notice Boxes (Lines 928-957)}

Three themed boxes using tcolorbox:

\begin{verbatim}
% Notice box (blue theme)
\newtcolorbox{notice}[1]{
    commonbox,
    colback=Azure1, colframe=DodgerBlue3,
    title={[Notice] #1}
}

% Definition box (green theme)
\newtcolorbox{definition}[1]{
    commonbox,
    colback=Honeydew1, colframe=ForestGreen,
    title={[Definition] #1}
}

% Conclusion box (red theme)
\newtcolorbox{conclusion}[1]{
    commonbox,
    colback=LavenderBlush1, colframe=Crimson,
    title={[Conclusion] #1}
}
\end{verbatim}

\subsection{Acknowledgement Environment (Lines 962-970)}

\begin{verbatim}
\newenvironment{acknowledgement}{
    \clearpage
    \addcontentsline{toc}{section}{Acknowledgement}
    \section*{Acknowledgement}
    \zihao{-4}\songti\linespread{1.5}\selectfont
}{
    \par
}
\end{verbatim}

% =========================================================
% Chapter 11: Module 17 - List Customization
% =========================================================
\newpage
\section{Module 17: Professional List Customization}

\subsection{Custom TikZ Bullet Markers (Lines 981-1005)}

\begin{verbatim}
% Level 1: Solid square
\newcommand{\hust@item@square}{%
    \tikz[baseline=-0.6ex] 
    \node[fill=black, shape=rectangle, 
          inner sep=0pt, minimum size=0.5em] {};%
}

% Level 2: Gray circle
\newcommand{\hust@item@circle}{%
    \tikz[baseline=-0.6ex] 
    \node[fill=darkgray, shape=circle, 
          inner sep=0pt, minimum size=0.35em] {};%
}

% Level 3: Dash line
\newcommand{\hust@item@dash}{%
    \tikz[baseline=-0.6ex] 
    \draw[line width=0.8pt] (0,0) -- (0.5em,0);%
}
\end{verbatim}

\subsection{Itemize Configuration (Lines 1018-1031)}

\begin{verbatim}
\setlist[itemize,1]{
    label=\protect\hust@item@square,
    leftmargin=2.0em, labelsep=0.6em
}
\setlist[itemize,2]{
    label=\protect\hust@item@circle,
    leftmargin=1.8em, labelsep=0.6em
}
\setlist[itemize,3]{
    label=\protect\hust@item@dash,
    leftmargin=1.8em, labelsep=0.6em
}
\end{verbatim}

\subsection{Custom Enumerate Styles (Lines 1045-1070)}

\begin{verbatim}
% Parenthesized: (1), (2), (3)
\newlist{enumpar}{enumerate}{1}
\setlist[enumpar]{label=(\arabic*)}

% Half-parenthesized: 1), 2), 3)
\newlist{enumbra}{enumerate}{1}
\setlist[enumbra]{label=\arabic*)}

% Circled numbers
\newlist{enumcircled}{enumerate}{1}
\setlist[enumcircled]{
    label=\protect\hust@circled{\arabic*}
}
\end{verbatim}

% =========================================================
% Chapter 12: Module 18-19 - Utilities and Review System
% =========================================================
\newpage
\section{Module 18-19: Utilities and Review System}

\subsection{Directory Tree Display (Lines 1077-1078)}

\begin{verbatim}
\RequirePackage{dirtree}
\renewcommand{\DTstyle}{\ttfamily\small}
\end{verbatim}

\subsection{Todo/Fixme/Note Commands (Lines 1100-1139)}

Breakable note boxes:

\begin{verbatim}
\newcommand{\todo}[1]{%
    \begin{tcolorbox}[
        hust@todostyle,
        colback=orange!5,
        colframe=orange!80!black,
        title={TODO (Pending)}
    ]
    #1
    \end{tcolorbox}%
}
\end{verbatim}

\subsection{PDF Metadata Injection (Lines 1142-1149)}

\begin{verbatim}
\AtBeginDocument{
    \hypersetup{
        pdftitle={\@title},
        pdfauthor={\@author},
        pdfsubject={HUST Report/Thesis},
        pdfcreator={HUST Report LaTeX Class}
    }
}
\end{verbatim}

\subsection{Review System (Lines 1166-1284)}

Using \texttt{pgfkeys} for option parsing:

\begin{verbatim}
\pgfkeys{
    /hustreview/.is family, /hustreview,
    style/.store in  = \hust@review@style,
    color/.store in  = \hust@review@color,
    side/.store in   = \hust@review@side,
    offset/.store in = \hust@review@offset
}

\NewDocumentCommand{\review}{ O{} m +m }{%
    \pgfkeys{/hustreview, #1}%
    % Render mark based on style (underline/wave/etc.)
    % Draw side note using TikZ overlay
}
\end{verbatim}

% =========================================================
% Chapter 13: Module 20 - Tutorial Documentation
% =========================================================
\newpage
\section{Module 20: Tutorial Documentation Environments}

\subsection{democode: Side-by-Side Display (Lines 1327-1351)}

\begin{verbatim}
\newtcblisting{democode}[1]{
    enhanced, breakable,
    title={#1},
    listing engine=listings,
    sidebyside,
    sidebyside align=top,
    lefthand ratio=0.5,
    bicolor,
    colbacklower=hust@demo@resultbg
}
\end{verbatim}

\subsection{demovert: Vertical Display (Lines 1356-1377)}

\begin{verbatim}
\newtcblisting{demovert}[1]{
    enhanced, breakable,
    listing engine=listings,
    listing and text,  % Code above, result below
    bicolor,
    colbacklower=hust@demo@resultbg
}
\end{verbatim}

\subsection{API Reference Card (Lines 1400-1418)}

\begin{verbatim}
\newcommand{\apibox}[3]{%
    \begin{tcolorbox}[
        colback=hust@api@bg,
        colframe=hust@api@frame,
        title={\textbackslash #1}
    ]
    \textbf{Parameters:} \texttt{#2} \\
    \textbf{Description:} #3
    \end{tcolorbox}
}
\end{verbatim}

\subsection{Inline Code Helpers (Lines 1438-1440)}

\begin{verbatim}
\newcommand{\cmd}[1]{\texttt{\textbackslash #1}}
\newcommand{\env}[1]{\texttt{#1}}
\newcommand{\opt}[1]{\texttt{[#1]}}
\end{verbatim}

% =========================================================
% Summary
% =========================================================
\newpage
\section*{Summary}
\addcontentsline{toc}{section}{Summary}

This document has provided a comprehensive analysis of the \texttt{hustreport.cls} class file, covering all 20 functional modules:

\begin{enumerate}
    \item Version declaration and class options
    \item Package loading organized by function
    \item User information interface
    \item Smart font configuration with fallback chain
    \item Conditional page layout (report vs. thesis)
    \item Section title and TOC styling
    \item Float optimization and caption formatting
    \item PDF properties and hyperlinks
    \item Special environments (abstract, bibliography, appendix)
    \item Cover and task book generation with TikZ
    \item Advanced text-figure layouts
    \item Professional code blocks with syntax highlighting
    \item English abstract environment
    \item Algorithm environment localization
    \item Themed notice/definition/conclusion boxes
    \item Acknowledgement environment
    \item Custom list markers with TikZ
    \item Utility features (dirtree, todo, metadata)
    \item Advanced review and annotation system
    \item Tutorial documentation environments
\end{enumerate}

\vspace{2cm}
\begin{center}
    \rule{0.4\textwidth}{0.5pt}\\[5pt]
    {\small Document generated with hustreport.cls v3.0}\\
    {\small \today}
\end{center}

\end{document}
